\section*{Chapter \ref{ch:b1:spherical-astronomy} --- Spherical Astronomy}

\begin{solution}{exo:b1:spherical-astronomy:1}
The \omterm{def:b1:spherical-astronomy:horizontal}{altitude} of the pole equals the \omterm{def:b1:spherical-astronomy:geographic}{latitude}: \ang{52}. A star never sets if
$\delta > \ang{90} - \ang{52} = \ang{38}$, and never rises if
$\delta < \ang{-38}$.
\end{solution}

\begin{solution}{exo:b1:spherical-astronomy:2}
$\qty{6.75}{h}\times\ang{15} = \ang{101.25}$. $H = \mathrm{LST} - \alpha =
\qty{10}{h} - \qty{6}{h}\,\qty{45}{min} = \qty{3}{h}\,\qty{15}{min}$, that is
\ang{48.75}. The \omterm{def:b1:spherical-astronomy:equatorial}{hour angle} is positive: the star is west of the meridian,
which it crossed three and a quarter sidereal hours ago.
\end{solution}

\begin{solution}{exo:b1:spherical-astronomy:3}
$h = \ang{90} - \varphi + \delta_\odot$: \ang{73.44} at the June solstice,
\ang{50} at the equinoxes, \ang{26.56} at the December solstice.
\end{solution}

\begin{solution}{exo:b1:spherical-astronomy:4}
1900: divisible by 100 but not by 400, not a leap year. 2000: divisible by
400, a leap year. 2100: not a leap year. The Julian year exceeds the \omterm{def:b1:spherical-astronomy:tropical-year}{tropical
year} by $365.25 - 365.2422 = 0.0078$ days, \num{0.78} days per century; the
Gregorian year by $0.0003$ days, \num{0.03} days per century.
\end{solution}

\begin{solution}{exo:b1:spherical-astronomy:5}
$\cos\Delta = \sin^2\ang{45} + \cos^2\ang{45}\cos\ang{90} = 0.5$, so
$\Delta = \ang{60}$ and the distance is $6371\times\pi/3 = \qty{6672}{km}$.
Along the parallel, a \omterm{def:b1:spherical-astronomy:spherical-triangle}{small circle} of radius $R\cos\ang{45}$, the distance is
$6371\times\cos\ang{45}\times\pi/2 = \qty{7076}{km}$. The \omterm{def:b1:spherical-astronomy:spherical-triangle}{great circle} is
shorter by about \qty{404}{km}.
\end{solution}

\begin{solution}{exo:b1:spherical-astronomy:6}
$\sin h = \sin\ang{48}\sin\ang{20} + \cos\ang{48}\cos\ang{20}\cos\ang{30}
= 0.2542 + 0.5446 = 0.7987$, $h = \ang{53.0}$.
$\cos\mathrm{A} = (\sin\ang{20} - \sin\ang{48}\sin\ang{53.0})/(\cos\ang{48}
\cos\ang{53.0}) = (0.3420 - 0.5935)/0.4027 = -0.6246$, $\arccos = \ang{128.7}$;
the star is west of the meridian ($H > 0$), so $\mathrm{A} = \ang{360} -
\ang{128.7} = \ang{231.3}$: toward the south-west, \ang{53} above the horizon.
\end{solution}

\begin{solution}{exo:b1:spherical-astronomy:7}
$\cos H_0 = -\tan\ang{60}\tan\ang{23.44} = -1.7321\times0.4335 = -0.7510$,
$H_0 = \ang{138.67}$, day length $2\times138.67/15 = \qty{18.49}{h}$, about
\qty{18}{h}~\qty{29}{min}, as the \ang{60} N curve shows at the June solstice.
\end{solution}

\begin{solution}{exo:b1:spherical-astronomy:8}
The passage moves earlier by \qty{3.93}{min} a day, so by $30\times3.93 =
\qty{118}{min}$ in 30 days: the star crosses at about 20:02. To cross at
noon it must move \qty{10}{h} = \qty{600}{min} earlier, which takes
$600/3.93 \approx 153$ days.
\end{solution}

\begin{solution}{exo:b1:spherical-astronomy:9}
(a) At an equinox the noon \omterm{def:b1:spherical-astronomy:horizontal}{altitude} is \ang{45}: $I = 1361\times\sin\ang{45}
= \qty{962}{W.m^{-2}}$. (b) At the June solstice $h = \ang{68.44}$ and
$(a/r)^2 = (149.598/152.100)^2 = 0.9674$:
$I = 1361\times0.9674\times\sin\ang{68.44} = \qty{1224}{W.m^{-2}}$. The
higher Sun more than makes up for the larger distance.
\end{solution}

\begin{solution}{exo:b1:spherical-astronomy:10}
$\sin\delta_\odot = \sin\ang{23.44}\sin\ang{60} = 0.3445$,
$\delta_\odot = \ang{20.15}$. $\tan\alpha_\odot = \cos\ang{23.44}\tan\ang{60}
= 1.5890$, $\alpha_\odot = \ang{57.82}$ (same quadrant as $\lambda_\odot$).
The \omterm{def:b1:spherical-astronomy:equatorial}{right ascension} lags by \ang{2.18}, that is $2.18\times4 \approx
\qty{8.7}{min}$ of time. A Sun moving uniformly along the \omterm{def:b1:spherical-astronomy:ecliptic}{ecliptic} would
still move unevenly in \omterm{def:b1:spherical-astronomy:equatorial}{right ascension}, the quantity that sets the \omterm{def:b1:spherical-astronomy:equatorial}{hour
angle} and hence the sundial: behind the mean Sun between an equinox and the
next solstice, ahead between a solstice and the next equinox. This is the
obliquity term of the \omterm{def:b1:spherical-astronomy:solar-time}{equation of time}.
\end{solution}

\begin{solution}{exo:b1:spherical-astronomy:11}
At rising $h = 0$, so \cref{prop:b1:spherical-astronomy:transform} gives
$\sin\delta = \cos\varphi\cos\mathrm{A}_0$; at the June solstice
$\delta = \varepsilon$, so $\cos\mathrm{A}_0 = \sin\varepsilon/\cos\varphi$.
For $\mathrm{A}_0 = \ang{45}$: $\cos\varphi = \sin\ang{23.44}/\cos\ang{45} =
0.5626$, $\varphi = \ang{55.8}$ N. Above $\ang{90} - \varepsilon =
\ang{66.56}$, $\sin\varepsilon/\cos\varphi > 1$: there is no \omterm{def:b1:spherical-astronomy:horizontal}{azimuth} of
rising, because the Sun does not set on that day (midnight Sun).
\end{solution}

\begin{solution}{exo:b1:spherical-astronomy:12}
(a) A lune of angle $\theta$ is the fraction $\theta/2\pi$ of the sphere:
$\frac{\theta}{2\pi}\times4\pi R^2 = 2\theta R^2$. (b) The three pairs of
lunes have total area $2\times(2A + 2B + 2C)R^2$; they cover the sphere once,
except the triangle and its opposite, each covered three times, that is
twice more than needed: $4(A + B + C)R^2 = 4\pi R^2 + 4S$, where $S$ is the
triangle's area. Hence $S = R^2(A + B + C - \pi)$. (c) Three right angles:
$S = R^2(3\pi/2 - \pi) = \pi R^2/2$, one eighth of the sphere (an octant).
The triangle of the figure: $89.9 + 64.9 + 67.4 - 180 = \ang{42.2}$, that is
\num{0.736} radians: $S = 0.736\,R^2$, \qty{5.9}{\percent} of the sphere.
\end{solution}

\begin{solution}{pb:b1:spherical-astronomy:1}
\textbf{1.} Colatitudes $\ang{90} - (\ang{-24.63}) = \ang{114.63}$ and
$\ang{90} - \ang{51.48} = \ang{38.52}$; difference of \omterm{def:b1:spherical-astronomy:geographic}{longitude} \ang{70.40}.

\textbf{2.} $\cos\Delta = \cos\ang{114.63}\cos\ang{38.52} +
\sin\ang{114.63}\sin\ang{38.52}\cos\ang{70.40} = -0.3260 + 0.1899 =
-0.1362$.

\textbf{3.} $\Delta = \ang{97.83}$; distance $6371\times1.7074 =
\qty{10878}{km}$, about \qty{10900}{km}.

\textbf{4.}
\[
  \cos\mathrm{Az} = \frac{\cos\ang{38.52} - \cos\ang{114.63}\cos\ang{97.83}}
  {\sin\ang{114.63}\sin\ang{97.83}} = \frac{0.7823 - 0.0568}{0.9006} = 0.8056 ,
\]
$\mathrm{Az} = \ang{36.3}$, toward the north-east.

\textbf{5.} $\sin\mathrm{Az} = \sin\ang{70.40}\sin\ang{38.52}/\sin\ang{97.83}
= 0.5922$, which gives \ang{36.3} or \ang{143.7}; the positive cosine of
question 4 selects \ang{36.3}.

\textbf{6.} No. A constant heading crosses every meridian at the same angle
and spirals toward a pole; the \omterm{def:b1:spherical-astronomy:spherical-triangle}{great circle}'s heading changes continuously
(it is \ang{36.3} only at the start). The aircraft would fly a longer path
and arrive elsewhere unless it corrected its heading.

\textbf{7.} The \omterm{def:b1:spherical-astronomy:horizontal}{altitude} of the south \omterm{def:b1:spherical-astronomy:celestial-sphere}{celestial pole} equals the \omterm{def:b1:spherical-astronomy:geographic}{latitude}'s
magnitude: \ang{24.63}.

\textbf{8.} $h = \ang{90} - |\ang{-24.63} - (\ang{-16.72})| = \ang{82.09}$;
$\delta > \varphi$, so Sirius passes \ang{7.91} north of the \omterm{def:b1:spherical-astronomy:celestial-sphere}{zenith}.

\textbf{9.} A star never sets if $\delta < -(\ang{90} - \ang{24.63}) =
\ang{-65.37}$. Canopus, at \ang{-52.70}, is not circumpolar: it sets.

\textbf{10.} $\ang{90} - |\ang{-24.63} + \ang{52.70}| = \ang{61.93}$, south
of the \omterm{def:b1:spherical-astronomy:celestial-sphere}{zenith}.

\textbf{11.} $\cos H_0 = -\tan(\ang{-24.63})\tan(\ang{-16.72}) =
-(0.4584)(0.3004) = -0.1377$; $H_0 = \ang{97.92}$. Sirius stays up
$2\times97.92/15 = \qty{13.06}{h}$ of \omterm{def:b1:spherical-astronomy:sidereal-time}{sidereal time}, about \qty{13}{h}.

\textbf{12.} $\cos\mathrm{A}_0 = \sin(\ang{-16.72})/\cos\ang{24.63} =
-0.3165$: Sirius rises at the \omterm{def:b1:spherical-astronomy:horizontal}{azimuth} \ang{108.5}, south of east.

\textbf{13.} When $H = 0$, $\mathrm{LST} = \alpha = \qty{6}{h}~\qty{45}{min}~\qty{9}{s}$.

\textbf{14.} $\qty{86400}{s}\times365.2422/366.2422 = \qty{86164.09}{s} =
\qty{23}{h}~\qty{56}{min}~\qty{4}{s}$; the passage moves earlier by
\qty{235.9}{s}, \qty{3}{min}~\qty{56}{s}, each night.

\textbf{15.} We need $12 + n\times24/365.2422 = 6.7525 + 24$, so
$n = 18.7525/0.065710 = 285.4$ days. 285 days after 20 March is
30 December: Sirius crosses the meridian at local mean midnight in the night
of 30 to 31 December 2026.

\textbf{16.} Paranal is \ang{70.40} west of Greenwich: local mean time is
behind by $70.40/15 = \qty{4.693}{h}$, \qty{4}{h}~\qty{41.6}{min}.

\textbf{17.} \omterm{def:b1:spherical-astronomy:solar-time}{Apparent solar time} = \omterm{def:b1:spherical-astronomy:solar-time}{mean solar time} + \omterm{def:b1:spherical-astronomy:solar-time}{equation of time} =
$\qty{0}{h} - \qty{3}{min}$: the sundial reads about 23:57.

\textbf{18.} $7\times\qty{3.93}{min} \approx \qty{27.5}{min}$ earlier.

\textbf{19.} $+\ang{23.44}$ at the June solstice, $-\ang{23.44}$ at the
December solstice.

\textbf{20.} June: $\ang{90} - |\ang{-24.63} - \ang{23.44}| = \ang{41.93}$.
December: $\ang{90} - |\ang{-24.63} + \ang{23.44}| = \ang{88.81}$; as
$\delta_\odot = \ang{-23.44} > \varphi$, the Sun passes \ang{1.19} north of the
\omterm{def:b1:spherical-astronomy:celestial-sphere}{zenith}.

\textbf{21.} $I = 1361\times(1/0.9837)^2\times\sin\ang{88.81} = 1361\times
1.0334\times0.9998 = \qty{1406}{W.m^{-2}}$.

\textbf{22.} In June, the southern winter. $\cos H_0 = -\tan(\ang{-24.63})
\tan\ang{23.44} = 0.4584\times0.4335 = 0.1988$: the Sun sets at the \omterm{def:b1:spherical-astronomy:equatorial}{hour angle}
$H_0 = \ang{78.53}$.

\textbf{23.} $\cos\mathrm{A}_0 = \sin\ang{23.44}/\cos\ang{24.63} = 0.4376$:
$\mathrm{A}_0 = \ang{64.0}$, north of east.

\textbf{24.} Day length $2\times78.53/15 = \qty{10.47}{h}$; the longest night
of the year at Paranal lasts $24 - 10.47 = \qty{13.53}{h}$, that is
\textbf{\qty[text-series-to-math=true]{13}{h}~\qty[text-series-to-math=true]{32}{min}}.
\end{solution}
